\section{Terminal trees and exposure of first defects}
\label{exp:section}
\label{sec:exposure}

The signed terms of Definition~\ref{cell:histories} are deterministic
whole-cell transports.  Randomness in this section means integration against
coordinates for such a term.  In particular, none of the auxiliary omissions
below is an operation on the microscopic process.  We use the fixed Gaussian
bound of Lemma~\ref{hist:gaussian}, the contact charts of
Lemma~\ref{hist:pinning}, and the exceptional-parameter and crowding estimates
of Propositions~\ref{hist:irregular} and~\ref{hist:crowding}.
Throughout the geometric applications use
\begin{equation}
\label{exp:time-scale}
 \tau_\varepsilon=\frac14\varepsilon^{.95}.
\end{equation}
The existing boundary exclusion of width $\varepsilon^{.95}$ leaves no
cell or layer cut in an initial interval of length $2\tau_\varepsilon$.
The retained birth spacing $2\varepsilon^{.9}$ is stronger than needed
on this shorter interval.  With arc error $C_M\varepsilon^{.96}$ the tube
ratio is $O(\varepsilon^{.01})$, and every fixed polynomial in $p$ times
$\tau_\varepsilon$ is $o(\varepsilon^{.85})$.

\subsection{Two coordinate systems, used at different stages}

The \emph{original forest} uses the first pin required by the factor in the
exact expansion: for a singleton addition, this is its first backward
contact with any root of that factor.  All ordinary parent counts,
compact-parameter exclusions, and increasing-tree estimates below use
these original coordinates.  A parent's identity determines its factor;
there is no additional choice of a factor after the parent is chosen.

The final exposure estimate uses a second, \emph{global forest}.  Its
anchored pool includes every displayed path, including paths from other
tensor factors.  A floating translation is pinned at its first backward
contact with this entire pool.  With several floating groups, their first
mutual contacts are used as well.  Coordinate groups merge at these pins,
but forces remain those of the original whole-cell prescription.  Thus a
new contact used only for this coordinate construction is virtual.  All
original required-contact, isolation, and first-contact tests remain
indicators.  This change of coordinates is used only after an extensive
number of fresh terminal trees has been selected, when its full
$p^{Cp}$ multiplicity is affordable.

A singleton global pin still produces exactly one new label.  A genuine
isolated scattering has both the continuing parent and the new child as
outputs.  A virtual contact or force-free passage changes only the new
branch.  A formerly prescribed genuine encounter following an earlier
virtual global pin remains an actual encounter, but is an additional
contact in the global forest.

Call a contact of the free added ray with another factor \emph{premature}
if it occurs before its original pin in backward time.  In the presence of
such a contact the global pin may be \emph{rerouted}.  The next lemma is
proved in the original coordinates, so it incurs no extra factor choice.
Bad original pins are handled by the original exceptional-parameter
bounds, rather than being assigned an unjustified small angular measure.

\begin{lemma}[Premature contacts in original coordinates]
\label{exp:rerouting}
Outside the original compact-parameter, crowded-parent, and multi-addition
exceptions, the contribution with a specified set of $k$ singleton cells
having a premature contact has bound
\begin{equation}
\label{exp:reroute-bound}
 C_L^p\bigl(p^C\varepsilon^c\bigr)^k.
\end{equation}
Consequently an extensive number of premature contacts has arbitrarily
strong exponential suppression.  The remaining premature births may be
marked irregular before selecting terminal subtrees.  No global
reparenting multiplicity is included or needed in this estimate.
\end{lemma}

\begin{proof}
There is exactly one added label across all expansions in the cell under
consideration.  Let $\sigma$ be its original backward pin time, and $a$ its
original parent.  The original first-pin condition says that on its
prelaunch interval the new label is free and has no contact with any root
of its own factor.  The current-cell root paths of all other factors are
also independent of its free velocity $w$ and position.  These paths are
known after the root states at the later cell endpoint are fixed.

Write $v=v_a(\sigma)$ and $g=w-v$.  With $u>0$ denoting elapsed time from
$\sigma$ in the prelaunch direction, the new ray is
\[
 x_a(\sigma)+\varepsilon n+u(v+g).
\]
The original parent is separated by $\varepsilon^{.85}$ from every other
already anchored path at time $\sigma$, outside the original crowding
exception.  Hence, with polynomial speed bounds, a premature contact is
impossible for $u<\tau_\varepsilon$.  For fixed $\sigma,n,|g|$, the remaining
angular measure is bounded by a constant times spherical area.  Indeed,
$w\mapsto g$ is a translation, the Gaussian density in $w$ is bounded
above on the retained compact range, and the incoming-hemisphere
restriction and bounded flux weight preserve this upper bound.
Uniformity of the original angular law is not needed.  The proof
of Lemma~\ref{exp:tube}, with time run in this prelaunch direction, bounds
the contact measure against each fixed other-factor root curve by
$p^C\varepsilon^{.01}$.  The union over target lives costs only a polynomial
factor.  An endpoint contact is included in the same estimate.

Use the product Gaussian/normal reference measure in the
Cauchy--Schwarz parent comparison of Proposition~\ref{hist:basic}, retaining all
first-pin tests as indicators.  Fix the entire ordered time vector and
the discrete original parent choices before this argument.  Index the
remaining coordinate blocks in backward cell order, including an entire
multi-addition block whenever necessary, and let $\mathcal F_i$ contain
all fixed times and choices, the top coordinates, and the first $i$ blocks.
A complete path or prelaunch segment from a previously processed block is
then fixed.  The current cell's roots and the other factors' current-cell
root paths are $\mathcal F_{i-1}$-measurable.  There is no other new label
in a designated singleton cell.  Its absolute velocity and normal are
fresh coordinates of block $i$.

Define $A_i$ using only these paths and block $i$: its original parent is
separated at the pin, the paths on the tested prelaunch interval have the
required polynomial speed bound, the relative velocity has the retained
compact bounds, and the free new ray meets an already anchored path of
another factor before its original pin.  In particular we do not include
any test against a future pin or condition on the eventual set of marks.
The separation and exterior-speed requirements are
$\mathcal F_{i-1}$-measurable.  On an original good history the global
energy bound implies these local speed requirements; outside that domain
we may simply omit the history.  The angular calculation gives, uniformly
in all past values,
\[
 A_i\in\mathcal F_i,\qquad
 \mathbb E_{\mathrm{ref}}[\mathbf1_{A_i}\mid\mathcal F_{i-1}]
 \le\rho_p:=p^C\varepsilon^{.01}.
\]
This is an adapted assertion, not a separate conditional bound with every
other block held fixed.  For any fixed set $S$ of $k$ designated indices,
let $i$ be its largest element.  The product of all other designated
indicators is $\mathcal F_{i-1}$-measurable.  Conditional expectation removes
$\mathbf1_{A_i}$ at cost $\rho_p$.  Iterating gives
\[
 \mathbb E_{\mathrm{ref}}\prod_{i\in S}\mathbf1_{A_i}\le\rho_p^k.
\]
All original-domain indicators not used in $A_i$ can be dropped for this
upper bound, without changing the normalized reference kernels.  Finally
sum over possible fixed sets $S$; this is a union bound, not conditioning
on which indices turn out to be marked.
The square-flux comparison contributes $C_L^p$ and at worst halves the
positive exponent.  Choices of target lives cost $p^{Ck}$; choices of the
$k$ cells cost at most $2^p$ when summed.  This proves
\eqref{exp:reroute-bound}.  Since $p$ is a fixed power of $P$, an extensive
$k$ supplies the claimed suppression.
\end{proof}

\subsection{Increasing trees and a terminal antichain}

Endpoint-overlap pins, including compulsory labels in a top endpoint
polymer of a positive component witness, are included among irregular
vertices.  Such witnesses need only $h_0$ distinguished spatial anchors;
no isolation of those anchors at the final endpoint is used here.  A
large endpoint polymer is paid for by its endpoint-volume pins before the
regular-tree count.  If only a small number of endpoint participants remain,
they create only that many irregular vertices and additional roots when
cut out.  Thus a witness with one distinguished anchor does not quietly
assume that its entire final component is a singleton.

For a rooted plane binary forest $\mathcal T$, let $s(v)$ be the number of
internal vertices in the subtree rooted at an internal vertex $v$.  A
continuing-parent edge is an ordinary child edge in this definition.

\begin{lemma}[Terminal subtree selection]
\label{exp:terminal}
Fix $0<\alpha\le1$ and an exponential accuracy $A>0$, with the layer
count and polynomial size exponent $a>0$ fixed.  There are constants
$\theta=\theta(\alpha,A)>0$, $M<\infty$, $p_0<\infty$, and
$\varepsilon_0>0$, independent of $\varepsilon$, such that the following
statement holds uniformly for
\[
 0<\varepsilon<\varepsilon_0,\qquad p_0\le p\le P^a,\qquad h_0\le\theta p.
\]
In the absolute history integrals with $p$ labels and $h_0$ tops, the
following exclusions have total cost at most $e^{-Ap}$ times the
unrestricted exponential majorant, after choosing the compact and
irregularity cutoffs sufficiently strongly:
there are more than $\alpha p/8$ insertion vertices whose original descendant
subtree has more than $M$ vertices, or too many vertices influenced by an
irregular insertion.  Here $M$ is a fixed constant, chosen before
$\varepsilon\downarrow0$.

After these exclusions, if at least $\alpha p$ distinct new-particle lives
have an extra contact or an uncharged defect, there are at least $c_{\alpha,M}p$ pairwise disjoint
regular terminal trees, of size at most $M$, each containing a defect.  For a
genuine launch its selected tree contains both output branches; for a
passage launch it need contain only the new branch.  The selection can
additionally ensure that no outside primary pin in the global coordinates
attaches to a selected tree, and that the original and global genealogies
agree on every selected tree.
\end{lemma}

\begin{proof}
For a fixed forest shape with $n$ internal vertices, the number of increasing
numberings is
\begin{equation}
\label{exp:hook}
 \frac{n!}{\prod_{v\in\mathcal T}s(v)}.
\end{equation}
Indeed, assign the smallest available number to each tree root, distribute
the remaining numbers between its two child subtrees with the appropriate
binomial coefficient, and repeat.  The resulting factorials cancel at
every internal edge and give \eqref{exp:hook}; the same multinomial
calculation distributes numbers between forest components.  A plane binary
forest with $h_0$ roots has at most $C^{n+h_0}$ shapes.  Ordered times in an
interval of length $T$ have volume $T^n/n!$.  Dropping the layer and lifetime
restrictions only increases this volume.  Consequently, if $\alpha n/8$
vertices have $s(v)>M$, summation over their shapes and numberings has bound
\begin{equation}
\label{exp:large-subtrees}
 C_T^{n+h_0} M^{-\alpha n/8}.
\end{equation}
The Cauchy--Schwarz comparison with uniform parent counting in
Proposition~\ref{hist:basic} can replace the last exponent by a fixed positive
fraction of it, which still allows arbitrary exponential suppression by
increasing $M$.  Quantitatively, choose
$\theta\le\min\{1/4,\alpha/8\}$, so that $n=p-h_0\ge p/2$.
Let $\zeta>0$ be the fixed exponent left by the parent comparison.
Choose $M$ so large that
\[
 \frac{\zeta\alpha}{16}\log M
 \ge A+\log C_*+2,
\]
where $C_*^p$ bounds the shape and fixed-layer constants in the comparison.
Then the large-subtree exclusion is exponentially small uniformly for
every admissible $p$, without requiring $h_0/p$ to tend to zero.

By Lemma~\ref{exp:rerouting}, include the few remaining premature births
among the irregular vertices.  This count still uses the original forest;
no global-parent multiplicity has been introduced.  Irregular vertices can
be cut out before this count.  Alternatively order
the pins inside an irregular cell consistently with attachment to the
anchored pool and refine them to binary vertices.  The additional orderings
are bounded by the factorial in that cell's number of labels.  The
multi-addition estimates in Proposition~\ref{hist:irregular} absorb these factors.
Cutting $r$ vertices creates at most $2r$ extra roots, so the shape count
remains exponential.
Here the irregular count concerns added insertion vertices and the old
lives targeted by their primary pins, not all old lives that the added
group may subsequently meet.  A cell with $m$ added labels has at most
$m$ primary attachment edges, even if its motion later touches many more
old labels.  Those secondary contacts remain defects on the corresponding
regular candidate trees; they are not removed as $O(m)$ contaminated
participants.  Choose the allowed irregular fraction only after $M$
is fixed.

A vertex of a subtree with at most $M$ internal vertices has at most $M$
ancestors whose subtrees also have at most $M$ vertices.  Thus $r$ irregular
vertices contaminate at most $Mr$ candidate launch vertices.  Distinct
new-particle lives have distinct birth vertices.  Excluding the $h_0$ top
lives and the contaminated candidates leaves a fixed fraction of the
$\alpha p$ affected lives.  Their candidate subtrees are nested or disjoint;
each contains at most $M$ candidate birth vertices.  Choose an inclusion
maximal candidate, discard its at most $M$ nested candidates, and repeat.
This produces a disjoint family with at least a $1/M$ fraction of the
remaining candidates.  There is one additional mark needed before this greedy selection.  Form
the global pins only as combinatorial data on the current configuration.
For each rerouted singleton birth mark its new global parent life, and for
each irregular multi-pin mark the lives to which its pins attach.  There
are at most a constant times the number of irregular participants such
marks.  Count a target at its actual attachment segment; a life persists
through a cut along the same continuation path, so no new mark is created
merely by crossing a layer.  An original bounded subtree containing a marked life is among at
most $M$ candidate ancestors of that life, with another constant to include
the life itself.  Remove these candidates also.  For example, take $h_0\le\alpha p/8$
and choose the irregular participant bound so that all child and target
marks together eliminate at most $\alpha p/8$ bounded candidates; a bound
$r\le c\alpha p/(M+1)$ with a sufficiently small absolute $c$ suffices.
Together with the $\alpha p/8$ large candidates, this leaves at least
$\alpha p/2$ affected candidates.  Greedy selection therefore retains at
least $\alpha p/(2(M+1))$ disjoint trees, up to integer rounding.  Now an outside global pin cannot attach to a
selected tree: an unchanged singleton attachment would already be its
original descendant, a rerouted attachment has a marked target, and a
multi-pin has a marked target.  No vertex within a selected tree is
rerouted.  Thus both genealogies agree there.  This target marking is
necessary; marking only the rerouted child's original ancestors would not
exclude attachment to an unrelated hidden parent.

The marks and the selected antichain can depend on the phase coordinates.
No probability estimate is obtained by conditioning on this selection.
In the later exposure bound we sum over its possible labels and choices at
cost $p^{Cp}$, paid by an extensive geometric gain.  Replacing a selected
passage subtree by its new-child part preserves disjointness and its
designated affected life.  Choose the irregular fraction after $M$, as
specified above, and increase the requested accuracy of its finitely many
exclusions by a fixed constant to sum them.  Finally take
$p_0\ge 8(M+1)/\alpha$ to absorb integer rounding; for example the number
of selected trees is then at least $\alpha p/(4(M+1))$.
All these choices are fixed before $\varepsilon\downarrow0$.
\end{proof}

\subsection{Launch coordinates and a tube estimate}

Histories are ordered backward, but $t$ in the geometric formulas denotes
physical time and velocities are physical velocities.  At a regular singleton pin,
write $v$ for the already known parent velocity and $w$ for the free added
velocity.  The change
\begin{equation}
\label{exp:relative-chart}
 g=w-v
\end{equation}
has Jacobian one: before the first pin the parent trajectory is independent
of the new label.  In insertion order these changes form a triangular
transformation, including when the parent velocity has been changed by
previous genuine scatterings.  Later relative parameters on a selected tree
will always be fixed in laboratory axes.

The measure on an entering sphere, or on its time reverse, is
$|g\cdot n|\,d\sigma(n)\,dg$.  Equivalently, with $s=|g|$, impact radius
$\rho\in(0,1)$, and normalized Haar measure $dR$ on $SO(3)$, it is
\begin{equation}
\label{exp:rotation-measure}
 8\pi^2s^3\rho\,ds\,d\rho\,dR.
\end{equation}
This follows by writing $dg=s^2ds\,d\omega$ and the impact area as
$\rho\,d\rho\,d\varphi$.  It uses no inverse of the scattering angle.  On
the retained compact parameter sets, $s$ and both genuine output increments
relative to the original parent are bounded above and away from zero.
Their body-frame vectors depend on $(s,\rho)$; common rotation multiplies
both by $R$.  Each nonzero vector separately has a uniform spherical
marginal.  The two output vectors need not be independent.

Galilean covariance gives a useful affine representation.  Fix all relative
parameters after a launch, its launch time $t_0$, and all its radial
parameters.  Up to the first defect, every descendant of output $i$ has
position
\begin{equation}
\label{exp:affine}
 x_a(t)=x_0+(t-t_0)v+(t-t_0)R a_i+G_a(t)+e_a(t),
 \qquad |e_a(t)|\le C_M\varepsilon^{.96}.
\end{equation}
Here $G_a$ is independent of $R$ and is Lipschitz with the retained speed
bound.  At a later birth the two velocities equal the current parent
velocity plus increments depending only on that birth's relative data.
Induction proves \eqref{exp:affine} for the broken rays.  Each exact
scattering arc has duration at most $\varepsilon^{.96}$ and bounded speed;
its replacement by an instantaneous broken ray changes positions by at
most a constant times that duration.  There are at most $M$ such arcs.
The displacement $\varepsilon n$ at an insertion is absorbed in the same
error.  This also proves the estimate during a later arc.

For two descendants diverging at a genuine vertex, subtract their
representations at that vertex.  The relevant coefficient is the nonzero
outgoing relative velocity, whose direction is again spherical under the
common rotation.  If they diverge at a passage, use the nonzero incoming
relative velocity.  Fixing parameters strictly below the divergence leaves
these directions available for integration.

\begin{lemma}[A ray with an independent exterior]
\label{exp:tube}
Let $\omega$ have normalized surface measure on $\mathbb S^2$, let
$0<r_0\le r\le r_1$, and let $H:[\tau,T]\to\mathbb R^3$ be independent of
$\omega$ and Lipschitz with constant at most $V$.  If $0<\delta\le\tau r_0/10$,
then
\begin{equation}
\label{exp:tube-bound}
 \int_{\mathbb S^2}\mathbf1_{\{\exists t\in[\tau,T]:
 |tr\omega-H(t)|\le\delta\}}\,d\omega
 \le C\frac{(V+r_1)\delta}{r_0^2\tau}.
\end{equation}
For the tree consequence, fix $M\in\N$ and $a>0$, let $p\le P^a$,
and retain the compact scattering, duration, birth-spacing, and boundary
gates of Proposition~\ref{hist:irregular}. Test at most polynomially many
exterior lifetime segments, independent of the tested launch rotation and
with polynomial Lipschitz speed. Each such segment must either be tested
only at elapsed times at least $\tau_\varepsilon$, or admit such an
independent extension back to the launch whose point there is at distance
greater than $\varepsilon^{.85}$ from the launch parent. The unchanged
parent of a passage is also allowed if it is free during the initial
$2\tau_\varepsilon$ interval. Prescribed scattering and passage encounters
are exempt. Under these hypotheses, a regular tree of at most $M$ vertices
has a first unprescribed internal or tested exterior contact with parameter
measure at most $p^C\varepsilon^c$, for some $c>0$, in the unconditioned
bounded-density reference measures described above.
\end{lemma}

\begin{proof}
On $[a,2a]\cap[\tau,T]$, take a mesh of spacing
$\delta/(2(V+r_1))$.  A hit implies that a mesh point satisfies the same
inequality with $2\delta$ in place of $\delta$.  There are at most
$C(1+a(V+r_1)/\delta)$ points.  At a fixed time the allowed directions are
empty or contained in a spherical cap with area at most
$C\delta^2/(r_0^2a^2)$.  The mesh union therefore has measure at most
$C(V+r_1)\delta/(r_0^2a)$.  Sum this geometric bound over
$a=2^k\tau$ to obtain \eqref{exp:tube-bound}.

Apply the estimate to \eqref{exp:affine}, with
$\delta=C_M\varepsilon^{.96}$ enlarged to include the interaction radius,
and $\tau=\tau_\varepsilon=\varepsilon^{.95}/4$.  On the earlier side of a physical
launch time use elapsed time $a=t_0-t$ and replace the spherical direction
by its negative; on the prelaunch side use $a=t-t_0$.  The quotient $\delta/\tau$ is
$O(\varepsilon^{.01})$.  Polynomial speed and finite choices of descendant
rays cost a factor $p^C$.
A lifetime segment tested only after $\tau_\varepsilon$ needs no
launch-separation condition: extend its independent comparison curve to
$[\tau_\varepsilon,T]$ without increasing its Lipschitz constant and apply
\eqref{exp:tube-bound}. The following initial-interval argument concerns
only segments that are tested before $\tau_\varepsilon$.
The launch parent is separated from all other
anchored paths by $\varepsilon^{.85}$, except for the already excluded
crowded-parent choices.  Since $\tau p^C=o(\varepsilon^{.85})$, no hit with these other exterior
paths occurs at smaller elapsed time.  For a passage tree the unchanged
launch parent remains visible and is not separated from the new child at
launch.  Their prescribed passage is exempt.  Separation from all other
paths and the absence of another birth or a cut in the initial
$2\tau_\varepsilon$ interval leave that parent free there.  The no-cut
assertion follows from \eqref{exp:time-scale} and the original boundary
gate; consequently no component from another cell can be installed during
this interval.  The nonzero relative straight ray cannot return after its
prescribed passage.  For later times that parent
is included among the exterior curves in \eqref{exp:tube-bound}.
For internal contacts, use the last
common divergence and use the retained exclusion of other births within
$2\varepsilon^{.9}>2\tau_\varepsilon$ of it.  The initial scattering itself is prescribed and finishes before this
interval ends.  The next free outgoing relative ray cannot return to its
partner without a later change of velocity.  A union over the at most
$M^2$ relevant divergence/ray pairs completes the proof.
\end{proof}

\subsection{The joint measure and chronological exposure}

We state the measure argument explicitly because a Gaussian upper bound
alone does not imply independence after conditioning.

For a product of independent input functions $u_1,\ldots,u_p$, set
\begin{equation}
\label{exp:source-factor}
 a_{i,k}=\sup_{|x-k|\le1,\,v}e^{\beta_0|v|^2}|u_i(x,v)|,
 \qquad
 \mathcal A_p=\prod_{i=1}^p\sum_{k\in\mathbb Z^3}a_{i,k}
 =\prod_{i=1}^p\|u_i\|_{\mathrm{Bol},\beta_0}.
\end{equation}
Any scalar small factor multiplying these inputs is left outside the
estimate.  Write $V_i^{\rm in}$ for the velocity at label $i$'s unique
independent source and define the input kinetic budget
\[
 \mathcal K=\frac12\sum_{i=1}^p|V_i^{\rm in}|^2=S/2.
\]
The once-counted energy inequality of Lemma~\ref{hist:gaussian} implies
polynomial speed bounds on the shell $\mathcal K\le Rp$, including the
current, departed, and not-yet-injected coordinate states.  In the next
proposition the phrase ``source factor'' always means the explicit number
$\mathcal A_p$ in \eqref{exp:source-factor}.

\begin{proposition}[Causal exposure for mixed prescriptions]
\label{exp:exposure}
Fix an admissible mixed whole-cell prescription with $p\le P^a$ labels and
a total energy budget $\mathcal K\le Rp$.  Fix a collection of $k$ disjoint
regular terminal trees as in Lemma~\ref{exp:terminal}.  Sum the absolute
activity-normalized chart integrals over the event that every selected tree has a
first defect.  The result is bounded by
\begin{equation}
\label{exp:exposure-bound}
 \mathcal A_p\,(pJ)^{C p}\bigl(p^C\varepsilon^c\bigr)^{k/2},
\end{equation}
The factor $(pJ)^{Cp}$ includes cell-segment choices.  Constants may
depend on the fixed layer count, compact cutoffs, $M$, and $R$.
\end{proposition}

\begin{proof}
\emph{The unconditioned charts.}
Fix the cells, force partitions, lifetimes, and labels of the translation
pins.  Inactive labels may be extended freely.  The augmented map from
independent input coordinates to the current/departed coordinates preserves
volume by Lemma~\ref{hist:gaussian}.  At a cell, keep the end coordinates of
all its roots anchored and decompose the added coordinates into velocities,
internal relative positions, and one common translation for each floating
group.  Grow the groups backward using the common first-contact convention.
Until a merger, each group's motion is independent of the other group's
common translation.  If the contact is at $t$ between labels $i$ and $j$,
its relative translation is
\begin{equation}
\label{exp:pin-jacobian}
 r=\varepsilon n+X_i(t)-X_j(t),\qquad
 |\det D_{(t,n)}r|
 =\varepsilon^2|n\cdot(V_i(t)-V_j(t))|.
\end{equation}
Internal coordinates in this formula are held fixed.  Once a relative
translation is fixed, regard its groups as a single coordinate group,
retaining the prescribed forces, including absence of force across a
virtual contact.  Ordering the internal merger edges before their enclosing
edges makes the full change of variables triangular.  Its Jacobian is the
product of \eqref{exp:pin-jacobian} and the endpoint-overlap volume factors.
No derivative of a many-particle scattering map is introduced.

Earliest-contact and no-prior-hit tests are kept as indicators.  Fix a
deterministic ordering of candidate contact pairs and resolve simultaneous
first contacts by that ordering; retain the resulting strata rather than
asserting that all ties are null.  The area formula applied to the
translation map in \eqref{exp:pin-jacobian} gives zero spatial-translation
measure to its nontransverse images.  Endpoint overlaps are represented
separately by their volume pins.  For each discrete merger forest and tie
stratum the inverse reconstruction is unique outside the nontransverse
images: the pin data recover its relative translations, and the prescribed
flows recover the coordinates.  There are at most $(pJ)^{Cp}$
merger/partner choices.  Singleton cells have one ordinary parent choice,
in the original coordinates.  In the global coordinates there can be
additional choices even for an irregular reroute; they are included in
the present $(pJ)^{Cp}$ bound.  This coarse multiplicity is used only in the present
geometric-gain estimate.  The factors $\varepsilon^2$ cancel the activity
normalization.  Multi-pin internal volumes and endpoint factors retain the
bounds of Lemma~\ref{hist:pinning}.

For regular singleton vertices belonging to the selected trees make
\eqref{exp:relative-chart}, followed by \eqref{exp:rotation-measure}.
Keep the absolute terminal velocities of every other floating component
among the known coordinates $\xi$; do not replace them by relative
velocities depending on a future anchor.  The parent state is known before that vertex,
so these substitutions remain triangular across cells.  Partition the
source positions into summability boxes.  The energy bound restricts all
velocities and displacements to polynomial ranges; source boxes contribute the weights $\prod_i a_{i,k_i}$ of
\eqref{exp:source-factor}.  Every integrated top position is within a
polynomial displacement of its own source box, giving a polynomial volume
per top.  The sum of the source-box weights is exactly $\mathcal A_p$.  On each box
choice, bound every remaining flux and compact radial density by a
polynomial in $p$, and enlarge coordinate domains to product ranges.
Let $d|\mathfrak I_{\boldsymbol k}|$ be the absolute chart measure of
the activity-normalized integral of Section~\ref{sec:history-estimates},
restricted to these source boxes.
It has the domination
\begin{equation}
\label{exp:product-domination}
 d|\mathfrak I_{\boldsymbol k}|\le
 \Bigl(\prod_i a_{i,k_i}\Bigr)(pJ)^{Cp}\,d\nu_0(\xi)
                    \prod_{\ell=1}^k d\nu_\ell(\theta_\ell).
\end{equation}
The reference measures are finite product measures.  Put the entire
ordered \emph{global-chart} time vector in $\xi$, including all selected launch times, internal
birth times, removals, and lifetime cuts.  It is fixed before exposure:
a visible parent must be removed at a known time, even while the output
velocities of its genuine launch remain hidden.  The block $\theta_\ell$
contains the relative velocities, radial impact parameters, and rotations
in tree $\ell$.  The variable $\xi$ contains all other chart coordinates
and the fixed genuine/passage types.  At a rerouted vertex, the original
required contact time is derived from the reconstructed flow and retained
only through its indicator; it is not an additional independent time.
Original and global pin times coincide on the selected trees.  The tube estimate is uniform in the
fixed time vector, including all retained spacing restrictions.  Normalize their
polynomially bounded total masses and absorb them in $(pJ)^{Cp}$.
Equation~\eqref{exp:product-domination} is applied to the entire joint
integral.  We never divide by a conditional mass of an admissible sector.
In particular, a rare sector is not asserted to have bounded normalized
conditional density.

\emph{Visible, dormant, and private states.}
Fix a list specifying the order of exposures, the one or two selected trees
revealed at each exposure, and the partner lives or internal divergences
and cell/ray segments involved.  There are at most $(pJ)^{Cp}$ such
discrete choices.  Each selected tree is revealed once.  The list does not fix
the exposure times; these will be recovered successively from the blocks
it reveals.  We first describe the state and its updates between exposures.

Within a cell the state has three parts:
anchored visible labels with known phase points; selected trees that have
already launched and remain hidden; and components of added labels that
have not yet merged with any rooted component.  Call the last components
\emph{dormant}.  Every dormant component is omitted from the bath, whether
or not it contains a selected label.  In particular, a nonselected
singleton or a multi-label addition is not represented in the bath by an
absolute prelaunch path whose later anchor has yet to be determined.

For a dormant component $C$ with known parameters, choose one terminal
position as origin.  Its terminal velocities, internal endpoint offsets,
and already specified internal merger pins determine relative paths
$q_i^C(t)$ and velocities $v_i^C(t)$, up to one common translation.  Use the
original force partition restricted to $C$ to compute these paths; a
virtual internal merger changes the coordinate component, not the force
partition.  These are functions of $\xi$ and already revealed blocks only.
A selected regular addition is a singleton in its cell.  Its as yet
unknown relative velocity is kept in its hidden block and is used only
when its known parent state is available at launch.  Internal additions
on an already hidden selected tree are handled privately in that tree,
using its block.  They are not part of the visible bath.  Its private paths
follow the scheduled isolated binary scatterings and free passages, with
all later relative parameters fixed in laboratory axes.  On the enlarged
product domain these paths continue through unlisted hypothetical contacts
without adding forces.  They therefore retain the affine representation
\eqref{exp:affine}.  Agreement with the prescribed whole-cell dynamics is
required only up to the tree's first defect on the retained domain.

Read the fixed global merger times in backward order.  If two dormant
components merge, the contact normal and time fix their relative
translation by \eqref{exp:pin-jacobian}.  Install that relative translation
inside their union and continue the restricted dynamics of the union,
retaining one free common translation.  It remains dormant.  No absolute
position or force from it is used in the visible evolution.  Its internal
mergers precede its first rooted merger in this construction.

Suppose next that $C$ first merges with the rooted pool at fixed time $t$,
through its label $i$ and a visible label $j$.  With the normal oriented
from $j$ to $i$, set
\begin{equation}
\label{exp:dormant-install}
 a_C=x_j(t)+\varepsilon n-q_i^C(t).
\end{equation}
Install every label $\ell\in C$ at
$(a_C+q_\ell^C(t),v_\ell^C(t))$ and thereafter evolve the enlarged rooted
system with its original prescribed forces.  The common translation uses
only the visible anchor and the known relative state.  Endpoint rooting
uses the corresponding overlap displacement in place of $\varepsilon n$.
For a selected launch this rule first computes its parent point and
velocity independently of its hidden block.  A genuine launch then hides
both outputs, including the continuing parent; a passage launch hides only
the new child.  No velocity of an unrevealed selected output is needed to
continue the remaining visible system.

The selection in Lemma~\ref{exp:terminal} rules out an outside primary pin
whose first rooted target belongs to a selected hidden tree.  This includes
rerouted singleton targets, multi-pin targets, and endpoint targets.  An
unchanged singleton pin into a hidden branch is its original descendant
and uses that tree's private block.  The launch antichain therefore ensures
that the parent of a different selected tree is visible when that tree
launches.

At a listed exposure, reveal the indicated block or pair of blocks,
install the current private states, and resume the original prescribed
force systems.  A virtual contact reveals the block without changing a
force.  An internal defect reveals its tree and resumes its prescribed
internal system.  Future additions whose data have just been revealed remain
dormant until their own merger times; revealing a block does not insert
all its future labels immediately.  No coordinate change at the random
exposure time is made.  The independent variables remain the global
whole-cell coordinates in \eqref{exp:product-domination}.

At a cell boundary carry
the visible root states to the next cell, retain the private states of
continuing hidden lines, and initialize that cell's new additions as
dormant.  A line that terminates is removed according to its prescribed
lifetime.  These updates use the global time vector already fixed in
$\xi$, including the removal time of a parent hidden by a genuine launch.

\emph{Recovering the listed exposure times.}
Suppose the first $m-1$ entries have been replayed.  The visible evolution
now uses only $\xi$ and those revealed blocks.  To replay entry $m$, supply
the one or two blocks named in that entry and construct their private
paths from their launch states.  Starting after the preceding replayed
exposure, find the first eligible contact of the listed type on its listed
segments, in backward chronological order with a fixed order for ties.
At that time make the revelation
update just described.  If there is no such contact, the list contributes
zero.  The same tie order permits successive entries at the same physical
time.  In this search we neither compare with a contact involving another
still-hidden block nor test that such a block has had no earlier hit.
These restrictions from the original first-exposure event are dropped.

The replay retains a revealed tree's private path before revelation and
its resumed path afterward.  They join at the installed phase point.
Thus a later test may use the whole past of a previously revealed life,
even when that life was hidden at the later-tested tree's launch.  Both
this past and the replayed revelation time depend only on the previously
revealed blocks.  In particular, revelation is not a new spatial placement.

\emph{Definition on the product domain and agreement.}
The enlarged product ranges need not satisfy the original energy and
isolation tests.  Choose a polynomial speed limit $V_p$ strictly larger
than the actual speed bound on $\mathcal K\le Rp$.  Stop each visible
simulation at its first violation of that limit or at failure of a validity
condition for its known finite flow.  Apply the same known-data rule to a
dormant relative flow needed for an installation.  Keep stopped positions
constant afterward, with zero velocity for any subsequent launch.  If a
prescribed visible anchor is unavailable, the list contributes zero and
its replay ends.  These rules use only visible
and already revealed data.  In particular, neither future contact
compatibility nor a dormant or hidden no-prior-hit test is used to continue
the simulation.  Its visible portions are $V_p$-Lipschitz for every value
of the product coordinates.

This exterior stopping rule is not a bound obtained by inspecting a hidden
tree.  Its private scheduled paths instead inherit the known launch-parent
velocity and add at most $M$ compactly bounded relative increments and
isolated scattering arcs.  Their speeds therefore have a polynomial bound,
uniform in the hidden rotations.  The common Galilean drift cancels in the
arc-to-ray error, leaving the bound $C_M\varepsilon^{.96}$ in
\eqref{exp:affine}.  Unlisted contacts do not alter those paths.  Their
private pasts and the resumed visible portions thus give polynomially
Lipschitz exterior curves.  The affine representation and tube estimate remain
available on the enlarged domain, without asserting an energy inequality
for its artificial trajectories.

We verify agreement on the retained chart stratum with the specified
actual exposure list.  Before a component's first rooted merger it has
no contact with an anchored path.  An earlier contact with a selected
hidden path would itself be an earlier rooted merger, and its outside
hidden target would have been excluded.  Omitting a dormant component
therefore changes no rooted motion before its scheduled installation.
This remains true within an original force group, since finite range
makes the cross force zero before contact.  A dormant--dormant merger
recovers the relative states up to one common translation, and
\eqref{exp:dormant-install} recovers that translation at a rooted merger.

Before the next actual first defect, the hidden private continuations
are their original isolated binary paths.  Thus the next listed contact
found by the replay occurs at its actual time and installs its actual
phase points.  Uniqueness then gives agreement of the resumed visible
flow.  Induction over the list proves agreement at every revelation;
the energy bound ensures that the visible stop is not reached before
any tested exposure.  No bound on the number of internal bath collisions
is used.  Connectivity roots every added component by the earlier cell
boundary on this stratum.  A hidden--visible contact at a cut is an exposure
or an excluded cut irregularity, and the intact-polymer condition prevents
a hidden force from being transferred silently through a cut.  Hence the
same induction passes through all cells and macro cuts.

We have defined curves on the entire product domain and proved that they
contain every hit event from the retained stratum.  Outside that stratum
they need only have the stated dependence on revealed blocks and the
uniform speed bounds.  No hidden coordinate has been conditioned on or
renormalized in this construction.

\emph{First-hit blocks.}
A hit of two hidden trees reveals them both and uses one estimate;
every other exposure reveals one tree.  Thus the fixed list has at least
$k/2$ entries.  Before its next test the bath and launch anchors depend
only on $\xi$ and previously revealed blocks.  Fixing the nonangular
parameters of the next block gives \eqref{exp:affine} against an
independent exterior.

For a one-tree test of $B$, the initial-interval hypothesis of
Lemma~\ref{exp:tube} also needs verification after earlier exposures.
Fix $\xi$ and the previously revealed blocks.  Their replayed paths,
including the private pasts retained at their revelations, use no parameter
of $B$.  Before its own launch the continuing parent of $B$ remains
visible.  On the retained domain these curves agree with the original
paths; on the enlarged domain their speed bounds were supplied above.

The launch cell of $B$ has exactly one addition across all factors,
and its boundary gate excludes a cut in the initial
$2\tau_\varepsilon$ interval.  Thus every exterior life tested before
elapsed time $\tau_\varepsilon$ was already displayed at $B$'s launch;
an earlier tree's revelation during this interval merely makes its
independent private past available.  Extend the tested life back to
that launch using the paths just constructed.  Except for the unchanged
parent of a passage, retain the gate that this extension starts farther
than $\varepsilon^{.85}$ from $B$'s launch parent, and declare the
enlarged early test empty if the gate fails.  The gate depends only on
$\xi$ and previously revealed blocks, and it holds on the selected
good domain by separated original parents and the agreement induction.
The initial offset is $O(\varepsilon)$ and
$V_p\tau_\varepsilon=o(\varepsilon^{.85})$, so the gate rules out an
early hit.

For the unchanged passage parent, first consider the retained domain.
At launch it is separated from every other old displayed path by
$\varepsilon^{.85}$.  The original polynomial speed bound and the
absence of another birth or cut in the initial $2\tau_\varepsilon$
interval preclude a third-path force there.  After the prescribed passage
the nonzero free relative ray cannot return to that parent.  Thus an
original first defect against the unchanged parent cannot occur at elapsed
time below $\tau_\varepsilon$.  On the product domain we consequently make
this parent test only at elapsed times at least $\tau_\varepsilon$;
no freeness of its artificial stopped path is asserted.  The tube estimate
then applies without an initial-interval gate.

Exterior segments born later
are tested only after $\tau_\varepsilon$; the same globally singleton
cell and boundary gate exclude an earlier installation.  For all tests
at elapsed times at least $\tau_\varepsilon$, apply the tube bound
without a launch-separation gate.  These extensions and gates are
independent of $B$'s tested rotation.

For a two-tree test, call $A$ the earlier-launched tree and $B$ the
later-launched tree in backward chronology; the retained spacing excludes
equal launch times.  Fix all of $A$'s block first and integrate $B$'s
launch rotation.  Both launch anchors were computed without either hidden
block.  Retain the gate that $B$'s launch parent is at least
$\rho=\varepsilon^{.85}$ from the already anchored paths of $A$ at that
launch.  The gate depends on $A$'s fixed block and the visible anchor of
$B$, not on $B$'s launch rotation.  If it fails, define this enlarged
pair test to be empty.  On the original selected good domain it holds by
parent separation.  During elapsed time $u<\tau=\tau_\varepsilon$ after
$B$'s launch, the initial offsets and displacements are bounded by
$C\varepsilon+Cp^C\tau=o(\rho)$.  The retained time spacing and boundary gate leave no
other birth or cut in this interval.  Thus no $A$--$B$ hit is possible there;
$B$'s own prescribed pair is not an $A$ path.  For $u\ge\tau$, apply
Lemma~\ref{exp:tube} to $B$ against the fixed paths of $A$.  The bound is
uniform in $A$'s block, so its subsequent integration preserves it.
Internal
defects use their divergence parameters as in that lemma.  The uniform
bound for each one- or two-tree test is
$\rho_p=p^C\varepsilon^c$.

For a fixed list, let $B_m$ denote the one- or two-tree block revealed
by its $m$th entry, and let $A_m$ require that its replay finds the listed
contact and satisfies the gates used in the corresponding test.  If the
replay has already failed, set $A_m=\varnothing$.  The construction and
the uniform one-test bound give
\[
 A_m=A_m(\xi,\theta_{B_1},\ldots,\theta_{B_m}),
 \qquad
 \int \mathbf1_{A_m}\,d\nu_{B_m}\le\rho_p,
\]
The blocks $B_m$ are disjoint, and the retained event for this list is
contained in $\bigcap_m A_m$.  The preceding no-hit requirements on
unexposed blocks have been dropped, without renormalizing their measures.
Integrate in reverse block order in
\eqref{exp:product-domination}.  Each last remaining restriction gives the
factor $\rho_p$, while all earlier restrictions are independent of that
block.  The resulting factor is $\rho_p^{\#\{B_m\}}$.
Summation over the guessed lists and the source-box weights, whose sum is
$\mathcal A_p$, proves
\eqref{exp:exposure-bound}.
\end{proof}

\subsection{Extensive defects}

\begin{proposition}[An extensive-defect gain]
\label{exp:defects}
Fix $0<\alpha\le1$, $A>0$, the layer count, and the polynomial size
exponent $a>0$.  There exist fixed constants $\theta=\theta(\alpha,A)>0$,
$p_0<\infty$, and sufficiently small $\varepsilon_0>0$ such that, after
the compact cutoffs are chosen, the following bound holds uniformly for
\[
 0<\varepsilon<\varepsilon_0,\qquad
 p_0\le p\le P^a,\qquad h_0\le\theta p.
\]
For histories with at least $\alpha p$ affected lives, their absolute
integral, divided by
the source factor $\mathcal A_p$ of \eqref{exp:source-factor}, is bounded
by the sum of the exceptional-parameter contributions bounded by $e^{-Ap}$
and a contribution bounded by
\begin{equation}
\label{exp:extensive}
 P^{C_Lp}\varepsilon^{c_Lp}.
\end{equation}
The constants $c_L>0$ and $C_L$ are fixed before $\varepsilon\downarrow0$.
The assertion includes recollisions, interference with a long-lived cluster,
and contacts with detached or subtracted trajectories.
\end{proposition}

\begin{proof}
We divide by the explicit source factor $\mathcal A_p$; if that factor
is zero the integral is zero.  First make the high-energy, multi-addition, endpoint, noncompact, long
scattering, close-time, and crowded-parent exclusions of
Propositions~\ref{hist:irregular} and~\ref{hist:crowding}.  Apply Lemma~\ref{exp:rerouting} to premature contacts in the original
coordinates.  Next choose $M$
and the irregular fraction in Lemma~\ref{exp:terminal}.  A positive fraction of the
affected lives then yields $k\ge c_{\alpha,M}p$ regular disjoint trees.
Choose $\theta$ and $p_0$ from Lemma~\ref{exp:terminal}, decreasing $\theta$
if the finite-participant transfer below requires a smaller affected
fraction.  Every exclusion is requested at a sufficiently increased fixed
accuracy so that their sum is at most $e^{-Ap}$.
Proposition~\ref{exp:exposure} gives \eqref{exp:extensive}.  Since
$p$ and $J$ are fixed powers of $P$, their logarithms are $O(\log P)=o(P)$,
and hence the expression in
\eqref{exp:extensive} is smaller than $e^{-Ap}$ for all sufficiently small
$\varepsilon$, uniformly in the indicated range of $p$: its logarithm
per label is at most $C_L\log P-c_L|\log\varepsilon|$.
For a fixed number $s$ of top roots, this applies to every
$p\ge m$ once
\[
 m\ge\max\{p_0,\lceil s/\theta\rceil\}.
\]
Thus it supplies a uniform tail in $p$, rather than a statement requiring
an $\varepsilon$-dependent limit of $s/p$.

The affected lives used by the algebraic centering argument transfer to this
geometric count with a fixed finite-participant loss.  An unchanged
algebraic extra contact marks its participating geometric lives.  If a
birth is rerouted to an earlier virtual global pin, mark its new child and
the old and new parents, at most three lives.  The later original required
contact is then a geometric defect of that child unless it is an already
charged irregularity.  A rerouted genuine contact cannot disappear by
calling it a different primary birth: each added label has exactly one
primary geometric pin.  This accounts for all discrepancies between the
two forests.  The charging statement of Proposition~\ref{cent:charges} may
therefore be used with a reduced positive fraction.

No estimate asserts that an already exposed bath has few collisions or no
bound cluster.  Such a bath is allowed in Lemma~\ref{exp:tube} through its
speed bound.  Its first influence on each fresh regular tree costs a
geometric factor; primary births into a small crowded bath are included in
the parent-choice exclusions.  This is why the argument treats interference
with clusters without deleting clusters from the actual dynamics.
\end{proof}

\subsection{A marked history of fixed size}

The extensive estimate is not needed, and does not directly apply, when
the number of displayed labels is fixed.  In particular, it does not cover
an array consisting only of several top roots.  We record the required
finite-size statement separately.

\begin{corollary}[Finite marked histories]
\label{exp:finite}
Fix the layer count and an integer $p_*$.  Sum the absolute integrals of the
original prescribed histories with at most $p_*$ labels, integrating all
top phase points, and suppose their independent source norms in
\eqref{exp:source-factor} are uniformly bounded.  The contribution with at
least one unprescribed contact, an incomplete isolated scattering at a
required cut, an interfering birth, or an endpoint-polymer defect tends to
zero as $\varepsilon\downarrow0$.  This statement includes $p=h_0$ and
contacts between distinct top roots.  The bounds are uniform over the
observation times and shortened endpoint cells in the fixed interval.
If instead a source has norm tending to zero, its finite-history integral
tends to zero by the same uniform finite-size majorant.
\end{corollary}

\begin{proof}
Work throughout in the original charts, without the coarse global-chart
multiplicity.  The ordinary parent/time bound of Proposition~\ref{hist:basic}
gives a finite constant depending on $p_*$ and $T$ after summing cell
locations.  Empty cells introduce no extra choice.  Its Gaussian bound
allows velocity truncation at a fixed large constant, with uniformly small
error.  The summability-box decomposition may be retained throughout;
its weights sum to \eqref{exp:source-factor}, so no spatial truncation
depending on the observation time is required.

Fix next $\eta>0$.  Delete births whose time is within $\eta$ of another
birth and parameters with incoming relative speed or a nonzero output
increment smaller than $\eta$.  The time error tends to zero with $\eta$
by the finite-dimensional simplex measure.  The parameter error tends to
zero by the flux-null statement of Proposition~\ref{pre:scattering} and
continuity from above of finite measure.  There is no claim of an
algebraic rate in $\eta$.  A scattering lasting more than
$\varepsilon^{.96}$ has vanishing flux measure.  Births within that distance
of a cell endpoint have total time measure at most
$C_{p_*}J\varepsilon^{.96}\to0$.  The multi-addition bounds give a vanishing
factor in a cell containing at least two additions, while an overlap pin
of an added label costs $\mu\varepsilon^3=\varepsilon$ in place of a
normalized flux.  A close pair of already integrated top roots instead
costs an unnormalized spatial volume $O(\varepsilon^3)$.  Endpoint polymers
are handled by a spatial spanning forest of these pins.  Thus no isolation
of a top array has been assumed in this deletion.

First test additional force-bearing encounters.  Before the first such
encounter every factor is a finite collection of isolated binary trees, with completed scattering arcs
replaced, at a position error
$\delta_\varepsilon=C_{p_*,\eta}\varepsilon^{.96}$, by their elastic broken
rays.  Purely virtual contacts do not change this reconstruction and are not
excluded at this stage.  The exact prescribed dynamics after the first
force-bearing event is unrestricted.
Only its volume and Gaussian majorants are used.  There are finitely many
ray pairs, divergence vertices, and birth/cut incidences to test.

Consider first a force-bearing contact between two families issuing from different top
roots.  Keep the internal collision data of both families fixed and float
one family by a common spatial translation $r$.  Before their first
contact, their motions are independent and translation equivariant.
The relative position is $r+G(t)$ with $G$ uniformly Lipschitz on the
compact parameter set.  The translations producing a hit within radius
$\delta$ have volume at most
\begin{equation}
\label{exp:spatial-tube}
 C\bigl(\delta^3+T\,\operatorname{Lip}(G)\delta^2\bigr).
\end{equation}
Indeed, sample the time interval at spacing
$\delta/(1+\operatorname{Lip}(G))$ and cover each sampled point by a ball
of radius $2\delta$.  Formula~\eqref{exp:spatial-tube} is uniform in the
locations of the source boxes.  Integrating their envelope weights proves
that root--root and cross-family contacts have vanishing integral.
The same calculation controls two top roots already close at an endpoint.
This argument uses the integrated top translations; it asserts no
pointwise smallness at fixed top positions.

For two rays within one family, use their last common divergence.  Its
nonzero relative output direction is spherical, even though the two
individual output velocities can be correlated.  Fix all relative data
below that divergence in laboratory axes.  Their relative position then
has the form used in Lemma~\ref{exp:tube}, with elapsed time at least
$\eta$ whenever another descendant has been born.  Before that time the
prescribed pair is a straight outgoing relative ray and cannot return.
Taking $\delta$ larger than
$\varepsilon+2\delta_\varepsilon$ gives a vanishing bound.  This proves
smallness of internal force-bearing recollisions without restricting what
happens after the first one.  Remove these vanishing physical-defect sets.
The complete remaining factor paths are then their prescribed binary
paths and free virtual rays throughout their lifetimes.  We may now test
all virtual contacts on these complete paths; they introduce no new
dynamical dependence.  Cross-family virtual contacts obey
\eqref{exp:spatial-tube}, and internal virtual contacts obey the same
divergence-direction argument.

The same two tests control a parent being within
$\varepsilon^{.85}$ of a third path at a scheduled birth: use the spatial
translation for different top families and the divergence direction for
one family.  The prescribed new child is exempt; all other divergences
are separated in time by the retained $\eta$.  Hence the near-parent
exclusions needed for finite singleton geometry have vanishing measure
without invoking the extensive crowding estimate.  For a premature
cross-factor contact of an added particle's full-cell prelaunch ray, use
its fresh incoming direction in the original chart, exactly as in
Lemma~\ref{exp:rerouting}; the comparison is with already anchored paths,
and no global reparenting is necessary.  The unchanged parent of a passage
is treated as in the proof of Lemma~\ref{exp:tube}.  Thus full-cell ghost
segments are included in the finite event list.

These estimates exhaust a first extra contact or interference of the
retained binary histories.  They also cover the short displacement and
cut errors, after the already removed duration and endpoint-time sets.
For fixed $p_*$, the number of combinatorial cases is finite and the time
sums are the original simplex integrals.  Take
$\varepsilon\downarrow0$ with the velocity and $\eta$ cutoffs fixed, then
let $\eta\downarrow0$ and the velocity cutoff tend to infinity.  The
resulting upper bounds depend on the observation time only through its
upper bound $T$; shortened cells only reduce their integration domains.
This proves uniformity.  Finally, replacing one source norm in the finite
majorant by a vanishing norm gives the last assertion directly.
\end{proof}
